\documentclass[12pt,a4paper]{article}

% ============================================================
% Paquetes
% ============================================================
\usepackage[utf8]{inputenc}
\usepackage[spanish]{babel}
\usepackage[margin=2.5cm]{geometry}
\usepackage{graphicx}
\usepackage{hyperref}
\usepackage{xcolor}
\usepackage{listings}
\usepackage{booktabs}
\usepackage{amsmath}
\usepackage{amssymb}
\usepackage{algorithm}
\usepackage{algpseudocode}
\usepackage{float}
\usepackage{caption}
\usepackage{tikz}
\usetikzlibrary{shapes,arrows.meta,positioning,fit}
\usepackage{enumitem}
\usepackage{longtable}

\hypersetup{
    colorlinks=true,
    linkcolor=blue,
    urlcolor=blue,
    citecolor=blue
}

\definecolor{codebg}{RGB}{245,245,245}
\definecolor{codekw}{RGB}{0,0,180}
\definecolor{codecomment}{RGB}{100,100,100}

\lstset{
    backgroundcolor=\color{codebg},
    basicstyle=\ttfamily\small,
    keywordstyle=\color{codekw}\bfseries,
    commentstyle=\color{codecomment}\itshape,
    breaklines=true,
    frame=single,
    numbers=left,
    numberstyle=\tiny,
    numbersep=8pt,
    tabsize=2,
    captionpos=b,
    language=C++
}

\title{
    \textbf{Reconstrucción Tridimensional de Órganos de Rana} \\
    \large Triangulación de Delaunay (Bowyer--Watson 3D) con Filtrado Alpha-Shape \\
    \vspace{0.3cm}
    \normalsize Laboratorio de Computación Gráfica --- UNSA-EPCC
}
\author{
    Rivas Huanca Diego Raul \\
    Alvarez Astete Jheeremy Manuel \\
    \small Escuela Profesional de Ciencia de la Computación \\
    \small Universidad Nacional de San Agustín de Arequipa
}
\date{\today}

\begin{document}

\maketitle

\begin{abstract}
Este informe documenta el diseño, implementación y validación de un
sistema de reconstrucción tridimensional de órganos de rana a partir de
máscaras volumétricas TIFF, utilizando triangulación de Delaunay en 3D
(algoritmo de Bowyer--Watson) combinada con un filtrado de tipo
alpha-shape para reconstruir superficies no convexas. El sistema
incluye un pipeline completo de procesamiento de datos (extracción de
superficie, sub-muestreo, cálculo de parámetros de reconstrucción) y
una aplicación interactiva de visualización en OpenGL/ImGui. Se valida
el algoritmo con casos sintéticos de topología conocida (esfera, cubo,
toro), confirmando su correctud mediante la característica de
Euler--Poincaré, y se aplica sobre los 17 órganos del conjunto de datos
real, documentando limitaciones, tiempos de ejecución y optimizaciones
de rendimiento (estructura de datos por hash, paralelización con
OpenMP).
\end{abstract}

\newpage

\tableofcontents
\newpage

% ============================================================
\section{Introducción y planteamiento del problema}
% ============================================================

La reconstrucción tridimensional a partir de datos volumétricos
segmentados es un problema central en visualización científica y
computación gráfica: dado un conjunto de máscaras binarias que indican
qué voxeles de un volumen pertenecen a una estructura de interés, se
busca generar una representación geométrica (malla de superficie)
navegable e interactiva de esa estructura.

En este laboratorio se trabaja sobre un conjunto de datos de 17
máscaras TIFF correspondientes a distintos órganos de una rana (sangre,
cerebro, corazón, hígado, músculo, esqueleto, bazo, entre otros), cada
una representando un volumen de $136 \times 470 \times 500$ voxeles. El
problema a resolver tiene dos componentes:

\begin{enumerate}
    \item \textbf{Geométrico}: dado un conjunto de puntos 3D extraídos
    de la superficie de un órgano, construir una malla triangulada que
    reconstruya fielmente su forma -- incluyendo concavidades reales
    del órgano, no solo su envolvente convexa.
    \item \textbf{De ingeniería de datos}: los volúmenes originales
    contienen entre $3{,}935$ (bazo) y $4{,}097{,}993$ (músculo)
    voxeles activos: es necesario un pipeline que extraiga, reduzca y
    prepare estos datos a un tamaño manejable para el algoritmo de
    reconstrucción, sin perder la posición relativa entre los distintos
    órganos del mismo espécimen.
\end{enumerate}

La dificultad central del componente geométrico es que un cierre
convexo (\emph{convex hull}) de la nube de puntos \textbf{no} reconstruye
correctamente órganos con concavidades -- que es el caso general en
anatomía. Esto motiva el uso de \emph{alpha-shapes}, una generalización
del cierre convexo que permite reconstruir formas no convexas a partir
de una triangulación de Delaunay.

% ============================================================
\section{Objetivos del proyecto}
% ============================================================

\subsection*{Objetivo general}
Implementar un sistema de reconstrucción 3D de órganos de rana a partir
de máscaras TIFF segmentadas, utilizando triangulación de Delaunay
(Bowyer--Watson en 3D) y visualización interactiva.

\subsection*{Objetivos específicos}
\begin{itemize}
    \item Implementar desde cero el algoritmo de Bowyer--Watson para
    triangulación de Delaunay en 3 dimensiones (tetraedros).
    \item Validar la correctud del algoritmo mediante propiedades
    matemáticas verificables (propiedad de Delaunay, característica de
    Euler--Poincaré) en casos sintéticos de topología conocida.
    \item Extender la triangulación con un filtrado alpha-shape para
    reconstruir superficies no convexas.
    \item Diseñar un pipeline de procesamiento que extraiga nubes de
    puntos representativas de la superficie de cada órgano a partir de
    las máscaras TIFF, preservando la posición relativa entre órganos.
    \item Implementar una aplicación de visualización interactiva
    (OpenGL + ImGui) que permita seleccionar, colorear y navegar los 17
    órganos reconstruidos.
    \item Medir y optimizar el rendimiento del sistema (estructuras de
    datos, paralelización con OpenMP, caché en disco de resultados ya
    calculados).
    \item Documentar las limitaciones del enfoque y posibles líneas de
    trabajo futuro.
\end{itemize}

% ============================================================
\section{Fundamento teórico}
% ============================================================

\subsection{Triangulación de Delaunay}

Dado un conjunto de puntos $P = \{p_1, \dots, p_n\} \subset
\mathbb{R}^3$, una triangulación de Delaunay es una descomposición del
casco convexo de $P$ en tetraedros tal que \textbf{ningún punto de $P$
se encuentra estrictamente dentro de la circunesfera de ningún
tetraedro} de la triangulación. Esta propiedad (la \emph{propiedad de
Delaunay} o \emph{propiedad del círculo/esfera vacía}) maximiza el
ángulo diedro mínimo de la triangulación resultante, evitando
tetraedros degenerados o extremadamente aplanados.

\subsection{Algoritmo de Bowyer--Watson}

El algoritmo de Bowyer--Watson \cite{bowyer1981,watson1981}
construye la triangulación de Delaunay de forma \textbf{incremental}:
se parte de un tetraedro (o triángulo, en 2D) suficientemente grande
como para contener a todos los puntos del conjunto (el
\emph{super-tetraedro}), y se insertan los puntos uno a uno. En cada
inserción:

\begin{enumerate}
    \item Se identifican todos los tetraedros cuya circunesfera
    contiene al nuevo punto (\emph{tetraedros inválidos}).
    \item Se eliminan esos tetraedros, dejando un hueco poliédrico.
    \item Se identifican las caras de frontera de ese hueco (caras que
    pertenecían a un solo tetraedro inválido).
    \item Se reconecta el hueco creando un nuevo tetraedro entre cada
    cara de frontera y el punto insertado.
\end{enumerate}

Al finalizar la inserción de todos los puntos, se descartan los
tetraedros que aún comparten algún vértice con el super-tetraedro
inicial.

\subsection{Circunesfera de un tetraedro}

Dado un tetraedro con vértices $A, B, C, D \in \mathbb{R}^3$, su
circunesfera es la única esfera que pasa por los 4 vértices. El centro
$O$ se obtiene resolviendo el sistema lineal $3\times3$ que surge de
imponer $\|O-A\|^2 = \|O-B\|^2 = \|O-C\|^2 = \|O-D\|^2$:

\begin{equation}
\begin{bmatrix}
(B-A) \\ (C-A) \\ (D-A)
\end{bmatrix} O
=
\frac{1}{2}
\begin{bmatrix}
\|B\|^2 - \|A\|^2 \\
\|C\|^2 - \|A\|^2 \\
\|D\|^2 - \|A\|^2
\end{bmatrix}
\end{equation}

Este sistema se resuelve mediante la regla de Cramer (determinantes
$3\times3$); el radio de la circunesfera es $r = \|O - A\|$.

\subsection{Alpha-shapes}

Un \emph{alpha-shape} \cite{edelsbrunner1994} es una generalización del
casco convexo parametrizada por un valor $\alpha$: dada la
triangulación de Delaunay completa de un conjunto de puntos, se
descartan los tetraedros cuya circunesfera tiene radio mayor a
$\alpha$, y se toma la frontera de los tetraedros restantes como
superficie reconstruida. Intuitivamente, $\alpha$ controla el tamaño
máximo de ``hueco'' que el algoritmo puede atravesar: valores grandes
de $\alpha$ tienden al casco convexo; valores pequeños fragmentan la
nube de puntos en componentes desconectadas.

\subsection{Característica de Euler--Poincaré}

Para una superficie triangulada, cerrada y orientable de género $g$
(número de ``agujeros''), se cumple:
\begin{equation}
V - E + F = 2 - 2g
\end{equation}
donde $V$, $E$, $F$ son la cantidad de vértices, aristas y caras. Para
una triangulación donde cada arista es compartida por exactamente 2
caras ($2E = 3F$), esto se reduce a:
\begin{equation}
F = 2V - 4(1-g)
\end{equation}
Para género 0 (topológicamente una esfera, sin agujeros): $F = 2V - 4$.
Para género 1 (un toro, con un agujero): $F = 2V$. Esta relación se usa
en este trabajo como \textbf{criterio de validación independiente} del
algoritmo (Sección~\ref{sec:validacion}).

% ============================================================
\section{Descripción del sistema desarrollado}
% ============================================================

El sistema se compone de dos partes que se comunican mediante archivos
en disco (nubes de puntos \texttt{.xyz} y un manifiesto
\texttt{manifest.tsv}):

\begin{enumerate}
    \item \textbf{Pipeline de datos (Python)}: procesa los TIFF
    originales, extrae la cáscara de voxeles de cada órgano, la
    sub-muestrea a una nube de puntos manejable, calcula estadísticas
    para sugerir un valor de $\alpha$, y exporta todo a disco.
    \item \textbf{Motor de reconstrucción y visualización (C++ /
    OpenGL)}: implementa el algoritmo de Bowyer--Watson, el filtrado
    alpha-shape, la construcción de mallas GPU, y una interfaz gráfica
    interactiva (ImGui) para seleccionar y navegar los órganos
    reconstruidos.
\end{enumerate}

Se desarrollaron tres ejecutables sobre el mismo núcleo de algoritmo:
uno de validación por consola (sin dependencias gráficas), uno de
validación visual con formas sintéticas y datos reales cargados a
mano, y el visor final del sapo completo (con y sin caché en disco).

% ============================================================
\section{Arquitectura general del proyecto}
% ============================================================

\begin{figure}[H]
\centering
\begin{tikzpicture}[
    box/.style={draw, rounded corners, minimum width=3.4cm, minimum height=1cm, align=center, font=\small},
    arr/.style={-{Latex}, thick}
]
\node[box, fill=blue!10] (tiff) {17 máscaras TIFF\\($136{\times}470{\times}500$)};
\node[box, fill=blue!10, right=1.2cm of tiff] (scripts) {Pipeline Python\\(scripts/)};
\node[box, fill=blue!10, right=1.2cm of scripts] (xyz) {.xyz + manifest\\(data/)};

\node[box, fill=green!10, below=1.2cm of scripts] (core) {\texttt{delaunay\_core}\\Bowyer--Watson 3D\\alpha-shape};
\node[box, fill=green!10, right=1.2cm of core] (cache) {\texttt{TetraCache}\\(.bin, opcional)};

\node[box, fill=orange!10, below=1.2cm of core] (gfx) {\texttt{app\_gfx\_core}\\GLFW + GLAD + ImGui};
\node[box, fill=orange!15, below=1cm of gfx] (apps) {Ejecutables:\\\texttt{delaunay\_test}, \texttt{sapo3d}, \texttt{sapo3d\_hd}};

\draw[arr] (tiff) -- (scripts);
\draw[arr] (scripts) -- (xyz);
\draw[arr] (xyz) -- (core);
\draw[arr] (core) -- (cache);
\draw[arr] (cache) -- (gfx);
\draw[arr] (core) -- (gfx);
\draw[arr] (gfx) -- (apps);
\end{tikzpicture}
\caption{Arquitectura general del sistema.}
\end{figure}

\subsection{Módulos principales}

\begin{description}
    \item[\texttt{src/delaunay/}] Núcleo del algoritmo, \textbf{sin
    dependencias gráficas}: \texttt{Vec3} (vector 3D en \texttt{double}
    para precisión numérica), \texttt{Delaunay3D}
    (Bowyer--Watson + alpha-shape + extracción de superficie),
    \texttt{IDelaunayStrategy} (patrón \emph{Strategy}),
    \texttt{PointCloudIO} / \texttt{ManifestIO} (lectores de texto
    plano), \texttt{TetraCache} (caché binario).
    \item[\texttt{src/core/}] Capa gráfica reutilizable:
    \texttt{Application} (ventana GLFW + contexto ImGui + loop
    principal), \texttt{Shader}, \texttt{Camera} (cámara orbital),
    \texttt{MeshBuilder} (construcción de buffers GPU).
    \item[\texttt{apps/}] Tres ejecutables: \texttt{synthetic\_test}
    (validación visual con formas sintéticas), \texttt{sapo3d}
    (reconstrucción completa), \texttt{sapo3d\_hd} (ídem, con caché
    persistente en disco).
    \item[\texttt{scripts/}] Pipeline de datos en Python:
    \texttt{tiff\_utils.py} (extracción de cáscara, sub-muestreo,
    estadísticas), \texttt{manifest\_utils.py} (lectura/escritura del
    manifiesto), \texttt{extract\_surface.py} / \texttt{extract\_all.py}
    (extracción de uno o todos los órganos), \texttt{report\_voxel\_counts.py}
    (reporte de referencia).
\end{description}

\subsection{Patrón de diseño: Strategy}

La interfaz \texttt{IDelaunayStrategy} desacopla el resto del sistema
de la implementación concreta del algoritmo de triangulación,
permitiendo -- sin modificar el código cliente -- intercambiar la
estrategia de reconstrucción (por ejemplo, para comparar variantes del
algoritmo o futuras optimizaciones):

\begin{lstlisting}[caption={Patrón Strategy sobre el algoritmo de triangulación}]
class IDelaunayStrategy {
public:
    virtual ~IDelaunayStrategy() = default;
    virtual std::vector<Tetrahedron> triangulate(
        const std::vector<Vec3>& points) = 0;
    virtual const char* name() const = 0;
};

class BowyerWatsonStrategy : public IDelaunayStrategy {
public:
    std::vector<Tetrahedron> triangulate(
            const std::vector<Vec3>& points) override {
        return BowyerWatson3D::triangulate(points);
    }
    const char* name() const override { return "Bowyer-Watson"; }
};
\end{lstlisting}

% ============================================================
\section{Metodología de implementación}
\label{sec:metodologia}
% ============================================================

\subsection{Bowyer--Watson en 3D}

El Algoritmo~\ref{alg:bowyerwatson} resume la implementación central.
La estructura de datos principal es un \texttt{Tetrahedron} con 4
índices hacia el arreglo de puntos, más el centro y radio al cuadrado
de su circunesfera (cacheados para no recalcularlos en cada test de
punto).

\begin{algorithm}[H]
\caption{Bowyer--Watson 3D}
\label{alg:bowyerwatson}
\begin{algorithmic}[1]
\Procedure{Triangulate}{$puntos$}
    \State $superTet \gets$ \Call{CrearSuperTetraedro}{$puntos$}
    \State $tets \gets \{superTet\}$
    \For{cada punto $p$ en $puntos$}
        \State $malos \gets \{t \in tets : p \text{ dentro de circunesfera}(t)\}$
        \State $caras \gets$ todas las caras de los tetraedros en $malos$
        \State $frontera \gets$ caras que aparecen una sola vez en $caras$
        \State $tets \gets (tets \setminus malos)$
        \For{cada cara $f$ en $frontera$}
            \State $tets \gets tets \cup \{\text{nuevoTetraedro}(f, p)\}$
        \EndFor
    \EndFor
    \State \Return $\{t \in tets : t \text{ no comparte vértice con } superTet\}$
\EndProcedure
\end{algorithmic}
\end{algorithm}

El super-tetraedro se construye tomando el centro y la diagonal del
\emph{bounding box} de los puntos de entrada, y generando 4 vértices en
configuración de tetraedro regular (vértices alternados de un cubo)
escalados a 5 veces esa diagonal suficientemente grande como para
garantizar que ningún punto de entrada quede fuera de él.

\subsection{Extracción de la superficie y alpha-shape}

\begin{lstlisting}[caption={Filtrado alpha-shape sobre la triangulación completa}]
std::vector<Tetrahedron> BowyerWatson3D::filterByAlpha(
        const std::vector<Tetrahedron>& tets, double alphaRadius) {
    double alphaSq = alphaRadius * alphaRadius;
    std::vector<Tetrahedron> result;
    for (const auto& t : tets) {
        if (t.radiusSq <= alphaSq) result.push_back(t);
    }
    return result;
}
\end{lstlisting}

La extracción de caras de frontera (\texttt{boundaryFaces}) identifica,
entre todas las caras de los tetraedros restantes, las que pertenecen a
un único tetraedro (las caras internas, compartidas por dos, se
descartan). Este paso se optimizó de $O(F^2)$ a $O(F)$ amortizado
mediante una tabla hash, ver Sección~\ref{sec:complejidad}.

\subsection{Pipeline de datos}

\begin{enumerate}
    \item \textbf{Extracción de cáscara}: de la máscara binaria sólida,
    se aplica erosión binaria y se resta del original, quedándose solo
    con los voxeles de borde. Reduce el conteo de puntos entre 5$\times$
    y 20$\times$ según el órgano antes de cualquier otro procesamiento.
    \item \textbf{Sub-muestreo por grilla 3D}: se particiona el espacio
    en celdas cúbicas y se toma un punto representante por celda
    ocupada, ajustando el tamaño de celda por búsqueda binaria hasta
    acercarse al conteo de puntos objetivo. Da una cobertura espacial
    más uniforme que un muestreo puramente aleatorio.
    \item \textbf{Centrado compartido}: todos los órganos se centran
    respecto al centro del volumen completo del TIFF (idéntico para los
    17 archivos, que comparten \emph{shape}), \textbf{no} respecto al
    centroide de cada órgano individual -- esto preserva la posición
    relativa entre órganos al combinarlos en la reconstrucción final.
    \item \textbf{Estadística de vecino más cercano}: se calcula la
    distancia mediana al vecino más cercano sobre la nube ya
    sub-muestreada, y se sugiere $\alpha \in [2.0, 3.0] \times$ esa
    mediana como punto de partida.
\end{enumerate}

% ============================================================
\section{Proceso de calibración del sistema}
% ============================================================

\subsection{Calibración del número de puntos objetivo}

Dado que el tamaño de la cáscara de voxeles varía en más de dos órdenes
de magnitud entre órganos ($1{,}385$ para el bazo hasta $822{,}268$
para el músculo), se calibró una función de escalado
$\text{target} = \text{clip}(K\sqrt{\text{cáscara}},\ 300,\ 3000)$, con
$K=16$ ajustado empíricamente a partir del caso del bazo (cáscara
$=1385 \rightarrow$ 600 puntos, resultado visualmente aceptable en las
primeras pruebas).

\subsection{Calibración de alpha}

El valor de $\alpha$ se calibró primero de forma \textbf{manual y
exhaustiva} sobre el caso sintético del toro (Sección~\ref{sec:validacion}),
donde el valor óptimo se pudo verificar de forma exacta contra la
predicción de Euler--Poincaré ($\alpha = 2.5$, sobre una nube con
distancia entre puntos vecinos $\approx 1.0$, es decir $\alpha \approx
2.5\times$ esa distancia). Ese factor ($2.0$--$3.0\times$ la distancia
mediana al vecino más cercano) se generalizó como heurística automática
para los 17 órganos reales, evitando así tener que ajustar el parámetro
a mano para cada uno.

\subsection{Validación cruzada con datos reales}

Se aplicó la heurística sobre el bazo (574 puntos, $\alpha_{auto}
\approx 2.5$) y el hígado (3039 puntos, $\alpha_{auto} \approx 9.6$),
confirmando visualmente formas reconocibles en ambos casos, con huecos
menores en zonas de baja densidad de puntos (extremidades finas)
documentado como limitación en la Sección~\ref{sec:limitaciones}.

% ============================================================
\section{Reconstrucción tridimensional}
% ============================================================

Una vez obtenidos los tetraedros filtrados por alpha-shape, la
reconstrucción de la malla de superficie sigue estos pasos:

\begin{enumerate}
    \item \texttt{boundaryFaces()} extrae las caras de frontera del
    conjunto de tetraedros filtrado.
    \item Por cada cara triangular, se calcula su normal (mediante
    producto cruz de dos de sus aristas) para \emph{flat shading}.
    \item Se construye un buffer de vértices intercalado
    (posición + normal, 6 \texttt{floats} por vértice, 3 vértices por
    cara, sin \emph{index buffer}) y se sube a la GPU
    (\texttt{MeshBuilder::buildSurfaceMesh}).
    \item El render usa un modelo de iluminación de Phong simplificado
    (ambiente + difusa + especular) en GLSL, con una cámara orbital
    controlada por arrastre de mouse y scroll.
\end{enumerate}

Para depuración también se construye una malla de líneas con las
aristas \'unicas de los tetraedros (sin filtrar por cara de frontera),
útil para inspeccionar visualmente la estructura interna de la
triangulación.

% ============================================================
\section{Procesamiento de la nube de puntos}
% ============================================================

El procesamiento de la nube de puntos ocurre enteramente en el lado
Python del pipeline (Sección~\ref{sec:metodologia}, pipeline de datos),
antes de que los datos lleguen al motor de triangulación en C++. Los
archivos intermedios (\texttt{.xyz}, texto plano con una fila
\texttt{"x y z"} por punto) se versionan de forma independiente de los
TIFF originales, lo que permite regenerar una nube de puntos con otro
nivel de detalle sin volver a tocar el algoritmo de reconstrucción.

El sistema de \emph{caché binario} (\texttt{TetraCache}) opera un nivel
más abajo: una vez que una nube de puntos fue triangulada, el resultado
completo (puntos + tetraedros, con su circunesfera ya calculada) se
persiste en disco (\texttt{data/cache/<órgano>\_<N puntos>.bin}),
evitando re-ejecutar Bowyer--Watson en corridas posteriores del
programa. El nombre del archivo de caché incluye la cantidad de puntos,
por lo que cambiar el nivel de detalle de un órgano no requiere
invalidar manualmente ningún caché existente.

% ============================================================
\section{Resultados obtenidos}
% ============================================================

\subsection{Validación con casos sintéticos}
\label{sec:resultados-sinteticos}

\begin{table}[H]
\centering
\caption{Resultados de validación con formas de topología conocida.}
\begin{tabular}{@{}lrrrl@{}}
\toprule
\textbf{Caso} & \textbf{Puntos} & \textbf{Tetraedros} & \textbf{Caras de frontera} & \textbf{Violaciones Delaunay} \\
\midrule
Esfera (genus 0)     & 40  & 123  & 76 (esperado 76, $2V{-}4$)  & 0 \\
Cubo aleatorio        & 60  & 285  & 46 (envolvente convexa)     & 0 \\
Toro sin $\alpha$ (genus 1) & 288 & 1524 & 332 (casco convexo, tapa el agujero) & 0 \\
Toro con $\alpha=2.5$  & 288 & 1128 & 576 (esperado 576, $2V$)    & 0 \\
\bottomrule
\end{tabular}
\end{table}

\subsection{Los 17 órganos}

\begin{table}[H]
\centering
\small
\caption{Conteo de voxeles originales y puntos finales usados por órgano.}
\begin{tabular}{@{}lrrrr@{}}
\toprule
\textbf{Órgano} & \textbf{Voxeles activos} & \textbf{Voxeles cáscara} & \textbf{Puntos usados} & \textbf{$\alpha_{auto}$} \\
\midrule
blood      & 35,492    & 26,035  & 2,492 & 5.00 \\
brain      & 19,172    & 5,374   & 1,200 & 3.54 \\
duodenum   & 38,224    & 12,182  & 1,741 & 3.54 \\
eye        & 53,231    & 9,024   & 1,525 & 3.54 \\
eyeRetna   & 37,146    & 13,317  & 1,832 & 3.54 \\
eyeWhite   & 16,085    & 5,165   & 1,158 & 3.54 \\
heart      & 37,196    & 8,648   & 1,561 & 3.54 \\
ileum      & 27,916    & 9,874   & 1,516 & 3.54 \\
kidney     & 39,245    & 12,199  & 1,753 & 3.54 \\
lIntestine & 51,660    & 9,634   & 1,625 & 3.54 \\
liver      & 205,245   & 42,951  & 3,039 & 5.00 \\
lung       & 47,454    & 13,210  & 1,782 & 3.54 \\
muscle     & 4,097,993 & 822,268 & 2,933 & 17.32 \\
nerve      & 24,314    & 16,649  & 1,987 & 3.54 \\
skeleton   & 440,640   & 198,340 & 3,015 & 8.29 \\
spleen     & 3,935     & 1,385   & 574   & 2.50 \\
stomach    & 134,323   & 67,602  & 2,975 & 5.59 \\
\midrule
\textbf{Total} & \textbf{5,309,271} & \textbf{1,273,857} & \textbf{32,708} & --- \\
\bottomrule
\end{tabular}
\end{table}

\subsection{Rendimiento}
\label{sec:complejidad}

\subsubsection{Complejidad algorítmica}

El algoritmo de Bowyer--Watson, tal como está implementado (sin
estructuras aceleradoras de tipo \emph{walk} o \emph{kd-tree}), tiene
complejidad $O(n^2)$ en el peor caso: por cada uno de los $n$ puntos
insertados, se recorre el conjunto de tetraedros existentes (que crece
proporcionalmente a $n$) para clasificarlos.

La extracción de caras de frontera (\texttt{boundaryFaces}) se
implementó inicialmente comparando cada cara contra todas las demás
($O(F^2)$, con $F=4\times|\text{tetraedros}|$), y se optimizó a
$O(F)$ amortizado usando una tabla hash sobre los 3 índices (ordenados
canónicamente) de cada cara:

\begin{table}[H]
\centering
\caption{Impacto de la optimización de \texttt{boundaryFaces()} (3000 puntos, 19{,}511 tetraedros, 78{,}044 caras).}
\begin{tabular}{@{}lr@{}}
\toprule
\textbf{Implementación} & \textbf{Tiempo} \\
\midrule
Comparación por pares, $O(F^2)$ & 2.20 s \\
Tabla hash, $O(F)$ amortizado & 0.005 s \\
\bottomrule
\end{tabular}
\end{table}

\subsubsection{Paralelización con OpenMP}

Dentro de cada inserción de Bowyer--Watson, la clasificación de
tetraedros (¿el punto cae dentro de su circunesfera?) es independiente
entre tetraedros, por lo que se paralelizó con
\texttt{\#pragma omp parallel for} (activado solo por encima de 500
tetraedros, para no pagar el overhead de creación de hilos en nubes
chicas):

\begin{lstlisting}[caption={Clasificación de tetraedros paralelizada con OpenMP}]
std::vector<char> isBad(tets.size(), 0);
#pragma omp parallel for schedule(static) if (numTets > 500)
for (long ti = 0; ti < numTets; ++ti) {
    if (inCircumsphere(tets[ti], p, allPoints, epsilonRel))
        isBad[ti] = 1;
}
\end{lstlisting}

\begin{table}[H]
\centering
\caption{Tiempos de triangulación secuencial vs. paralelo (OpenMP).}
\begin{tabular}{@{}lrrr@{}}
\toprule
\textbf{Puntos} & \textbf{Tetraedros} & \textbf{Secuencial} & \textbf{OpenMP (\underline{\hspace{1cm}} núcleos)} \\
\midrule
3000 & \textasciitilde19{,}500 & \underline{\hspace{1.5cm}} s & \underline{\hspace{1.5cm}} s \\
\bottomrule
\end{tabular}
\end{table}
\emph{Nota: Tabla con mediciones en la máquina de
desarrollo (multi-núcleo); el entorno usado durante el desarrollo
inicial del algoritmo solo contaba con un núcleo disponible.}

Se evaluó el uso de CUDA para acelerar aún más la triangulación, y se
descartó: Bowyer--Watson incremental es inherentemente secuencial
\emph{entre} inserciones (el estado de la malla tras insertar el punto
$i$ es precondición para insertar el punto $i+1$), por lo que no se
presta a paralelización masiva en GPU sin reemplazar el algoritmo
completo por un enfoque distinto (p. ej. \emph{divide and conquer}
paralelo). Además, la GPU del sistema ya se emplea activamente para el
render (shaders GLSL, buffers de vértices), que es la etapa que
efectivamente se beneficia de paralelismo masivo en este proyecto.

% ============================================================
\section{Validación del sistema y análisis de resultados}
\label{sec:validacion}
% ============================================================

\subsection{Propiedad de Delaunay}

Para cada tetraedro resultante, se verificó que ningún punto de la nube
de entrada (fuera de sus 4 vértices) cae estrictamente dentro de su
circunesfera. En los 4 casos sintéticos evaluados (esfera, cubo, toro
con y sin $\alpha$), el conteo de violaciones fue \textbf{0} en todos
los casos.

\subsection{Validación topológica mediante Euler--Poincaré}

El caso del toro es la prueba de validación más informativa del
proyecto: al ser una superficie de género 1 (con un agujero real), la
característica de Euler--Poincaré predice $F=2V$ caras en una
triangulación cerrada correcta. Con $V=288$ puntos, eso son 576 caras
esperadas. El filtrado con $\alpha=2.5$ produjo \textbf{exactamente
576} caras de frontera -- una coincidencia exacta que confirma que el
alpha-shape reconstruye la \emph{topología real} del objeto (agujero
incluido), y no solamente su casco convexo (que en el mismo caso dio
332 caras, un valor topológicamente incorrecto para un toro).

\subsection{Análisis de resultados sobre datos reales}

La aplicación sobre los 17 órganos reales confirma que el enfoque
escala razonablemente entre estructuras pequeñas y compactas (bazo) y
grandes/complejas (hígado, con lóbulos), produciendo reconstrucciones
reconocibles en ambos extremos. Sin embargo, se identificaron
limitaciones que se detallan a continuación.

\subsection{Limitaciones y dificultades}
\label{sec:limitaciones}

\begin{itemize}
    \item \textbf{Alpha-shape de radio único por órgano}: un solo valor
    de $\alpha$ no siempre reconstruye bien zonas anchas y finas del
    mismo órgano simultáneamente. En el bazo y el hígado aparecen
    huecos menores en extremidades/protuberancias finas, donde la
    distancia entre puntos vecinos crece localmente y el mismo
    $\alpha$ que conecta bien el cuerpo principal ya no alcanza ahí.
    Una alpha-shape \emph{adaptativa} (radio variable según densidad
    local) resolvería esto, pero excede el alcance de este laboratorio.
    \item \textbf{Sub-muestreo con tope fijo de puntos}: \texttt{muscle}
    (822K voxeles de cáscara) y \texttt{skeleton} (198K) se comprimen
    al mismo techo de $\sim$3000 puntos que órganos mucho más chicos
    como el bazo (1385) esto implica menos detalle relativo en las
    estructuras más grandes y complejas del espécimen.
    \item \textbf{Complejidad algorítmica sin aceleración espacial}: la
    implementación usa búsqueda lineal sobre los tetraedros existentes
    en cada inserción; es adecuada para las nubes de algunos miles de
    puntos usadas aquí, pero no escalaría bien a nubes de cientos de
    miles de puntos sin una estructura de aceleración espacial
    (p.~ej. \emph{walk} guiado por localidad, o un \emph{kd-tree}).
    \item \textbf{Supuesto de voxel isotrópico}: se asume que todas las
    máscaras comparten el mismo \emph{shape} de volumen y que los
    voxeles son isotrópicos (mismo espaciado en los 3 ejes) válido
    para este conjunto de datos puntual, pero no es un supuesto general
    para cualquier dataset médico (p.~ej. tomografías con distinto
    espaciado entre slices y dentro de cada slice).
\end{itemize}

% ============================================================
\section{Conclusiones y trabajo futuro}
% ============================================================

\subsection*{Conclusiones}

\begin{itemize}
    \item Se implementó y validó exitosamente el algoritmo de
    Bowyer--Watson para triangulación de Delaunay en 3 dimensiones,
    confirmando su correctud mediante dos criterios independientes: la
    propiedad de Delaunay (0 violaciones en todos los casos) y la
    característica de Euler--Poincaré (coincidencia exacta en el caso
    del toro, $F=576=2V$).
    \item El filtrado alpha-shape resultó necesario y efectivo para
    reconstruir formas no convexas -- sin él, el sistema reconstruiría
    únicamente el casco convexo de cada órgano, perdiendo toda
    concavidad anatómica real.
    \item El pipeline de datos (extracción de cáscara + sub-muestreo
    por grilla + centrado compartido) permitió procesar los 17 órganos
    de forma automática y consistente, preservando su posición relativa
    dentro del espécimen.
    \item Las optimizaciones de rendimiento (tabla hash para
    \texttt{boundaryFaces}, caché binario en disco, paralelización con
    OpenMP) redujeron significativamente el tiempo de uso interactivo
    del sistema, sin alterar la correctud de los resultados (verificado
    comparando resultados antes/después de cada optimización).
\end{itemize}

\subsection*{Trabajo futuro}

\begin{itemize}
    \item Alpha-shape adaptativa (radio variable por región, según
    densidad local de puntos) para reducir los huecos en extremidades
    finas.
    \item Estructura de aceleración espacial (kd-tree o \emph{walk})
    para bajar la complejidad práctica de la inserción de puntos,
    permitiendo subir el tope de puntos por órgano sin degradar el
    tiempo de respuesta interactivo.
    \item Suavizado de malla (p.~ej. Laplacian smoothing) como
    posproceso, para reducir el aspecto ``faceteado'' de las
    superficies reconstruidas con pocos puntos.
    \item Exploración de un enfoque de Delaunay paralelo por
    \emph{divide and conquer} si se buscara aprovechar GPU/CUDA para la
    etapa de triangulación (no solo para el render).
    \item Extensión del pipeline de datos para soportar voxeles
    anisotrópicos (espaciado distinto por eje), generalizando su uso a
    otros conjuntos de datos volumétricos médicos.
\end{itemize}

% ============================================================
\begin{thebibliography}{9}

\bibitem{bowyer1981}
Bowyer, A. (1981). \emph{Computing Dirichlet Tessellations}. The Computer Journal.

\bibitem{watson1981}
Watson, D. F. (1981). \emph{Computing the $n$-dimensional Delaunay Tessellation with Application to Voronoi Polytopes}. The Computer Journal.

\bibitem{edelsbrunner1994}
Edelsbrunner, H., \& Mücke, E. P. (1994). \emph{Three-Dimensional Alpha Shapes}. ACM Transactions on Graphics.

\end{thebibliography}

\end{document}